The environment uses the actual Accellera UVM 1.2 source library bundled with ModelSim Intel FPGA Starter 2020.1. A preliminary library probe executed factory construction, test phases and objections before the DUT environment was admitted. The completed DUT regression, rather than that probe, supplies the numerical evidence.

An \texttt{epoch\_item} carries case identity and the reset-abort selection. The sequence submits one interrupted epoch followed by four complete epochs through a UVM sequencer and driver. The driver controls PCM presentation, start/reset and independent readiness of the waveform and packet consumers through a virtual interface. Stimulus is scheduled away from the sampling edge. The passive monitor observes accepted transfers at the rising edge and sends separate analysis transactions to the scoreboard.

The scoreboard compares each accepted PCM token, padded source token, waveform output, canonical-bin record and exported packet byte with immutable admitted fixtures. It checks record identity and ordering as well as values. Final time-domain metrics are calculated independently by absolute output index. A second Python checker reads the recorded observations and reconciles them with the fixtures after the simulator exits.

The fixed corpus contains 1,024 input samples per completed epoch, 1,536 full-tail output samples, nine frames and 1,161 canonical-bin records. Thresholds are zero and $10^8$; each is exercised in dense and sparse packet mode. Three explicit seeds drive distinct reproducible stall schedules. For each seed, the first completed case is also the exact recovery check after the interrupted prefix. Recovery epochs must not be added a second time to the completed-epoch denominator.

\paragraph{Tool capability boundary.}
The Starter simulator reports restrictions on advanced verification features. This environment uses deterministic seeded UVM sequences and observed event counters. It does not claim solver-constrained randomization, executed covergroups, simulator code-coverage closure or a UVM-aware debugger. The library is compiled with \texttt{UVM\_NO\_DPI}; the disabled DPI-dependent facilities are not part of the acceptance contract. UVM 1.2 is identified explicitly, without a claim of IEEE 1800.2-2020 library conformance.
